\begin{viuabstract}{Abstract}
This thesis studies cooperative transport of heterogeneous payloads through
three subproblems: SP1 forms coalitions and assigns roles; SP2 approaches,
docks, transports, and replaces members; SP3 coordinates routes and
reservations across missions. The primary \emph{Cargo} mode models a supported
payload as a planar rigid body.

The proposal combines LP/MILP references, potential games, primal--dual
dynamics, integer closure, and contracts for resources, contact, motion,
recovery, and communication. Formal results cover quota regulation, local Cargo
stability under fixed contact and exact wrench realisation, termination of route
improvements, and the impossibility of guaranteeing global optimality when
disconnection hides locally indistinguishable instances with incompatible optima. Experiments use paired worlds, explicit seeds, and global
oracles identified as experimental ceilings.

SP1 regulates quotas under explicit conditions, separates heterogeneous
capacity, and reduces nominal false feasibility at the SP1--SP2 interface. SP2 exhibits an adverse
safety--progress trade-off and restores all recoverable certificates in the
static benchmark; the worst equilibrium of its repair game has no uniform
bound. In SP3, resource reservation reduces
deadlocks in the discrete catalogue; the congested pilot makes no delivery. Cargo
completes almost all missions without executing all prior theoretical
mechanisms as one chain.

The contribution is a verifiable architecture of contracts and local games,
together with a planar functional composition. It does not establish global
hybrid stability, social optimality, fully distributed coordination, or
industrial validity. Instrumented contact, hardware, and caging remain outside
the reported confirmatory evidence.

\keywords{Keywords}{multi-robot coordination,
  autonomous mobile robots, coalition formation,
  potential games, cooperative transport}
\end{viuabstract}

\clearpage
